\section{Ordinary rank-three points and automorphic descent}
\label{sec:ordinary-global-descent}

Let $\Omega$ be an algebraic closure of $\mathbf Q$, let $K$ be the marked imaginary quadratic field and write $G_F=\operatorname{Gal}(\Omega/F)$ for an embedded subfield. Fix the coefficient ring $\mathcal O$ at $\ell$, $\iota:\overline{\mathbf Q}_\ell\simeq\mathbf C$, the positive weight parameter $e$, the even positive modulus $N$, and the character $\psi$ of the good datum. Write $\epsilon$ for the $\ell$-adic cyclotomic character, $d$ for the character obtained from $\psi$ on the Galois group, and $\Psi=\epsilon^e d$, after embedding the coefficients. The character $\eta$ cuts out $K$. Fix the marked conjugation $c$, the place $B$ over $\ell$, and the element $g_1$ supplied by the good frame. For every finite $\mathcal O$-subalgebra $A\subset\overline{\mathbf Q}_\ell$ which is a discrete valuation ring and every marked ordinary point
\[
                R:G_{\mathbf Q}\longrightarrow
                 \mathrm{GL}_3(A)\rtimes\{1,-1\},
\]
with the sign acting by the framed inverse-transpose involution, let $\rho$ be its matrix representation on $G_K$ after extending the coefficients. Suppose that these matrices act irreducibly on every open subgroup of $G_{\mathbf Q}$ contained in $\ker\eta$. Section~\ref{sec:spectrum-open-irreducibility} proves this condition for the indicated type I points. The conclusion of the descent is that
\[
                       r=\rho\otimes\epsilon^{-1}
\]
is attached to a regular algebraic conjugate self-dual cuspidal representation over $K$ and satisfies the additional local assertion below. We give the reduction to the linked automorphy input over the auxiliary field and the cyclic base-change and attachment inputs in [1--6].

\subsection{Field, marked point and polarization}

Put $M=\ell N|\operatorname{disc}K|$. Any inertia group with nontrivial $\eta$ lies at a prime dividing $|\operatorname{disc}K|$. Indeed, at a prime away from the discriminant two embeddings of $\mathcal O_K$ which agree after reduction at a chosen prime of $\Omega$ agree in $\Omega$. This follows from the invertibility there of the discriminant pairing, or from the unramified local decomposition of $\mathcal O_K$. An inertia element acts identically on the residue field, so the original embedding and its translate agree. This gives the ramification condition at $M$, and $\ell N\mid M$. The parity assumption $\psi(-1)=(-1)^e$ and $\epsilon(c)=-1$ give $\Psi(c)=1$. Complex conjugation squares to the identity. Thus the condition at $B$, the ramification condition, and the divisibility and positivity conditions used in the link hold with the original $e,N$.

The imported frame theorem supplies an abelian CM field $L/\mathbf Q$, a frame $\Phi$, and a Galois link with $\Phi$ satisfying the two size and admissibility predicates and the seed condition of [1]. The link gives $G_L\subset\ker\eta$. For any allowed embedding of $K$, its image is the fixed field of $\ker\eta$ because the two embeddings of a quadratic Galois extension have the same image. Hence every element of that image is fixed by $G_L$ and belongs to $L$. This constructs $K\hookrightarrow L$ compatibly with the given embeddings in $\Omega$. It gives a finite abelian Galois extension $L/K$. Since $G_L$ is open in $G_{\mathbf Q}$ and lies in $\ker\eta$, $\rho|_{G_L}$ is irreducible. A nonzero scalar twist preserves the invariant subspaces. In particular $r|_{G_L}$ is irreducible.

Here is the precise use of automorphy over $L$. For every finite discrete valuation coefficient extension as above, the rank-three conclusion of the link applies to a normalized frame lift whose local inertia characters are $\Psi^{a_i}$, $(a_1,a_2,a_3)=(1,0,-1)$. It applies to any irreducible representation over $L$ whose matrices are the image of that lift evaluated on the map from $G_L$ to the frame group. Construction of the normalized lift from $R$ gives those matrices as $\rho|_{G_L}$, and its three inertia formulas give exactly these powers. The ring map $\mathcal O\to A$ is local by integrality. The rank-three twist built into the link is precisely $\epsilon_L^{-1}$. Consequently the conclusion over $L$ applies to $r|_{G_L}$. In [1] this link uses an ordinary rank-two seed, a rank-three density theorem for the unitary family, and propagation of the adjoint lift of the seed. These are global inputs even when one starts with an already irreducible marked point.

Let $J$ be the rank-three reversal matrix, so $J^2=1$. The normalized matrix component at $c$ is the identity, and its sign is $-1$. For $s\in G_K$ and $s'=csc^{-1}$ the equality $s'c=cs$ in the semidirect product therefore gives
\[
             \rho(s')=J\,{}^t\!\rho(s)^{-1}J,\qquad
             {}^t\!\rho(s)J\rho(s')=J .
\]
As $\epsilon(s')=\epsilon(s)$, twisting both factors gives the exact multiplier
\begin{equation}
 {}^t\!r(s)Jr(s')=\epsilon(s)^{-2}J
                 =\epsilon(s)^{1-3}J.                  \label{eq:ord-descent-pair}
\end{equation}
The multiplier on the real subfield $K^+$ is the same cyclotomic power. It is algebraic by the cyclotomic character construction, with local exponent $-2$ and trivial finite factor, and its value at each complex conjugation is $1$. The conjugation $c$ fixes $K^+$ because its image under every relevant real embedding is real.

Continuity after extending the coefficients can also be checked explicitly. Every element of $A$ is integral over $\mathbf Z_\ell$ and has norm at most one. Write $\ell=u\varpi^h$ for a unit $u\in A^\times$ and uniformizer $\varpi$. Then every element of $\mathfrak m_A^{hj}$ is divisible by $\ell^j$ and has norm at most $|\ell|_\ell^j$. Each entry of $R$ is locally constant modulo $\mathfrak m_A^{hj}$ by the adic continuity condition of the point. Taking $j\to\infty$ gives continuity of the entries of $\rho$. The inverse entries are continuous by applying the same argument at $s^{-1}$. Multiplication by the continuous scalar $\epsilon^{-1}$ gives continuity of $r$.

\subsection{Inertia away from \texorpdfstring{$\ell$}{l} and flags above it}

Fix a finite place $v\nmid\ell$ of $K$, a prolongation, and its decomposition and inertia groups $D_v,I_v$. The inertia group acts trivially on every $\ell^j$-th root of unity. Indeed reduction is injective on roots of unity of order prime to the residue characteristic, and inertia acts trivially on the reductions. Thus $\epsilon(I_v)=1$. The character $\Psi|_{G_K}$ is continuous and its restriction to $I_v$ has image in the finite image of $d$. If $s\in I_v$ and $\Psi(s)=1$, then $\eta(s)=1$ and the away-from-$\ell$ condition in the marked ordinary point gives $R(s)=1$. The same holds for $r(s)$ since $\epsilon(s)=1$.

The image $\Psi(I_v)$ is a cyclic group of roots of unity, of order $m$, say. The representation of $I_v$ factors through this group. Choose $u\in I_v$ mapping to a generator $\xi$. The operator $r(u)$ satisfies $X^m-1$, which is separable in characteristic zero. In an eigenbasis its eigenvalues are $\xi^{b_1},\xi^{b_2},\xi^{b_3}$ with $b_i\ge0$. For every inertia element $s$ it follows that
\[
       \det(X-r(s))=\prod_{i=1}^3(X-\Psi(s)^{b_i}).
\]
Each $\Psi^{b_i}$ extends continuously to $D_v$ by restriction of the global character. This is the full local characteristic-polynomial condition called \texttt{IneAb} in [2].

For $v\mid\ell$, choose a place of $\Omega$ over $v$. The frame conjugacy condition transports it to the marked place $B$. Conjugating the ordinary flag back supplies the upper triangular form of $\rho$ there. For a transporter with positive sign this uses matrix conjugation. For a transporter with negative sign it also uses $b\mapsto J\,{}^tb^{-1}J$. This involution takes an upper triangular matrix to an upper triangular matrix and sends its diagonal $(x_1,x_2,x_3)$ to $(x_3^{-1},x_2^{-1},x_1^{-1})$. Since $a_{4-i}=-a_i$, it preserves the prescribed three powers of $\Psi$ in their order. The ordinary characters of the successive one-dimensional quotients of $\rho$ therefore restrict on inertia to $\epsilon^{ea_i}d^{a_i}$. Twisting them all by $\epsilon^{-1}$ gives
\[
        \epsilon^{\,e-1}d,\qquad
        \epsilon^{-1},\qquad
        \epsilon^{-e-1}d^{-1}.
\]
Transport to any other prolongation over $v$ by a Galois element conjugates the flag and the decomposition group together. This proves the flag hypothesis for all prolongations in [2].

The descent of the local window uses a further property of the selected frame. Write $U_1(r,v)$ if at each prolongation there is a decomposition-group flag with diagonal characters $\kappa_i$ such that $\kappa_i\ne\kappa_j\epsilon$ for every $i,j$. Write $U_\infty(r,v)$ for the analogous condition with $\kappa_i^a\ne(\kappa_j\epsilon)^a$ for every integer $a\ge1$. The release proves for the even level frame that for each $w\mid\ell$ of $L$, with $v=w|_K$,
\[
             U_1(\rho,v)\ \Longrightarrow\
                     U_1(\rho|_{G_L},w).
\]
This preservation uses the linked noncoincidence theorem in [2]. A common scalar twist cancels from the two sides of every character equality, including their positive powers. Thus the same implication holds for $r$.

\subsection{Correcting the descent at a prime cyclic layer}

We detail the part of the global argument that identifies the specified representation downstairs. For a cyclic extension $E/F$ of prime degree in the solvable CM tower, let $r_F$ be the restriction of $r$ to $G_F$ and $r_E=r_F|_{G_E}$. Both are irreducible, since their groups contain $G_L$. Suppose that a regular algebraic cusp form $\Pi$ over $E$ is attached to $r_E$. At all but finitely many places, Frobenius classes at conjugate places over $F$ have conjugate matrices under $r_F$. Thus the attached Satake multisets of $\Pi$ at those places agree with their conjugates. The prime-degree cyclic descent assertion supplies a cusp form $\pi_0$ over $F$, weakly base changing to $\Pi$ and carrying the descended archimedean parameters [3,4]. Every complex embedding of $F$ extends to $E$, so $\pi_0$ inherits regular algebraicity. The attachment of a Galois representation downstairs requires an additional argument. The dual and the conjugate predicted by the polarization have weak base changes related by the predicted character, using the upstairs Frobenius data. Two descents of the same weak base change differ by a finite norm character, by the fiber assertion of the prime-degree theorem. In the release this norm character is converted into a Galois character on $G_F$ and extended to $G_{F^+}$. Dividing the previous polarizing character by this finite-order correction leaves its value constant on complex conjugations. Comparison of unramified Satake multisets and strong multiplicity one then polarize $\pi_0$ with this corrected character. The polarized Galois-attachment theorem supplies a continuous semisimple $r_0$ attached to $\pi_0$ [5].

The restriction of $r_0$ to $G_E$ is semisimple by normality and finite index. On restriction, its unramified parameters are those of $\Pi$, by the weak base-change identity. Uniqueness of the Satake multiset gives equality of the characteristic polynomials of $r_0|_{G_E}$ and $r_E$ on almost all Frobenius classes. The Frobenius isomorphism criterion (Chebotarev density and the characteristic-zero semisimple character criterion) gives
\[
                     T\,r_0(h)T^{-1}=r_F(h)
                         \qquad(h\in G_E).
\]
To see exactly what remains on $G_F$, put $r'_0=T r_0 T^{-1}$. For $g\in G_F$ both $r_F(g)$ and $r'_0(g)$ act by the same conjugation on $r_E(G_E)$. For $h\in G_E$ this follows by applying the equality to $ghg^{-1}\in G_E$. Hence $r'_0(g)^{-1}r_F(g)$ commutes with $r_E(G_E)$. Schur's lemma makes it a scalar $\delta(g)$ since the coefficients are algebraically closed and $r_E$ is irreducible. Multiplicativity of both representations gives $\delta(gg')=\delta(g)\delta(g')$. The scalar is one on $G_E$, so $\delta$ is a finite character of $\operatorname{Gal}(E/F)$. Use its corresponding finite-order Hecke character in the attachment convention of the release to twist $\pi_0$. Twisting preserves cuspidality and regular algebraicity and multiplies the attached representation by $\delta$. This yields an automorphic representation attached to $r_F$ itself.

Iteration is legitimate on this tower. The finite abelian group $\operatorname{Gal}(L/K)$ has a composition series with cyclic prime quotients. The intermediate fields contain $K$, and are CM. Since $L/\mathbf Q$ is abelian, complex conjugation induces a central automorphism of each intermediate field; its fixed field there is totally real, while its restriction to $K$ is nontrivial. Thus each intermediate field is quadratic over its totally real subfield. At every step the restriction of the original $r$ to the upper group is irreducible, and the original polarization restricts. The release proves this induction more generally for solvable CM extensions [3]. Applying it down to $K$ produces a regular algebraic cusp form attached to $r$. Equation~\eqref{eq:ord-descent-pair} then identifies its conjugate with its contragredient. In the attachment normalization the factor $\epsilon^{1-3}$ is precisely the prescribed duality twist. The release applies its comparison of dual Satake parameters and strong multiplicity one, using the trivial auxiliary Hecke character [5].

\subsection{The additional local assertion}

For a continuous representation, the assertion called \texttt{WitA} consists of two clauses [6]. For each $v\mid\ell$, $U_1(r,v)$ implies $U_\infty(r,v)$. For every regular algebraic cusp form $\pi$ attached to $r$, there exists a real $t$ such that, simultaneously for every finite place $v$ with $v\nmid\ell$ or $U_\infty(r,v)$, its local component satisfies the source condition $\operatorname{WinAt}(\pi,t,v)$. In this condition the support satisfies the toral condition $\operatorname{TorAt}(\pi,v)$, and for every character constituent $\xi$ of any such local cuspidal support one has, with $q_v$ the residue cardinality and $\varpi_v$ the chosen uniformizer,
\[
         q_v^{\,t-1/2}<|\xi(\varpi_v)|<q_v^{\,t+1/2}.
\]
The parameter $t$ may depend on $\pi$ but does not depend on $v$.

For the first clause take a place $w$ above $v$. The even-level preservation established above gives $U_1(r|_{G_L},w)$ from $U_1(r,v)$. The first clause over $L$ gives $U_\infty(r|_{G_L},w)$. The flag over $K$ descends this power inequality. Restricting a hypothetical equality of two powers at $v$ would give the equality at $w$. The multiset of diagonal characters of any two flags on the same restriction is equal, by equality of traces and linear independence of distinct characters in characteristic zero, so a change of the flag above $w$ cannot remove the equality. For the second clause use the upstairs form with the global parameter $t$ furnished by its local assertion. At a prime-degree step weak base change relates the two attached forms. The window comparison of [6] applies at every $v\nmid\ell$ using the inertia characteristic-polynomial condition, and at $v\mid\ell$ satisfying $U_\infty$ using the ordinary comparison at $\ell$. The same $t$ works at each such $v$. This is the second input beyond the descent of the cusp form. The linked density theorem over $L$, the polarized attachment used in the cyclic step, and this local window comparison are used as theorem inputs here.

\paragraph{Numbered sources.}
\begin{enumerate}
\item The release's \href{https://github.com/Deicyde/riemann-hypothesis/blob/40e56bab026552e089d7afcfeff08d2defb9c005/src/Development/AutomorphicArithmetic/AutomorphicVarieties/ModularityArgument/ThresholdArguments/DenseThresholdCases/GdEndOf.lean}{good frame and rank-three application} and \href{https://github.com/Deicyde/riemann-hypothesis/blob/40e56bab026552e089d7afcfeff08d2defb9c005/src/Development/AutomorphicArithmetic/AutomorphicVarieties/ModularityArgument/ThresholdArguments/DenseThresholdCases/AutEnd3Alt.lean}{rank-three endpoint proof} give the marked quantifiers and the choice of twist. The \href{https://github.com/Deicyde/riemann-hypothesis/blob/40e56bab026552e089d7afcfeff08d2defb9c005/src/Development/AutomorphicArithmetic/AutomorphicVarieties/ModularityArgument/ThresholdArguments/DenseThresholdCases/SymTop.lean}{linked density theorem} identifies the rank-two seed and unitary rank-three inputs.
\item The proofs of the \href{https://github.com/Deicyde/riemann-hypothesis/blob/40e56bab026552e089d7afcfeff08d2defb9c005/src/Development/AutomorphicArithmetic/AutomorphicVarieties/ModularityArgument/UpperTriangularPlaces/OrdinaryFlagThree.lean}{ordinary flag at each place}, the \href{https://github.com/Deicyde/riemann-hypothesis/blob/40e56bab026552e089d7afcfeff08d2defb9c005/src/Development/AutomorphicArithmetic/AutomorphicVarieties/ModularityArgument/AutomorphyLifting/AutomorphyDescent/InertiaRankTwo/LocalConditionRankThreeLinkOfInertiaOrdinary.lean}{away-from-$\ell$ inertia condition}, and the \href{https://github.com/Deicyde/riemann-hypothesis/blob/40e56bab026552e089d7afcfeff08d2defb9c005/src/Development/AutomorphicArithmetic/AutomorphicVarieties/ModularityArgument/ThresholdArguments/DenseThresholdCases/DnEnds.lean}{even-level preservation for the linked field} give the separate local hypotheses. The \href{https://github.com/Deicyde/riemann-hypothesis/blob/40e56bab026552e089d7afcfeff08d2defb9c005/src/Development/AutomorphicArithmetic/AutomorphicVarieties/CurveDVRSpecialization/OrthogonalPairingFromIntertwiner.lean}{matrix pairing and coefficient continuity lemmas} prove the corresponding assertions before descent.
\item The release's \href{https://github.com/Deicyde/riemann-hypothesis/blob/40e56bab026552e089d7afcfeff08d2defb9c005/src/Development/AutomorphicArithmetic/AutomorphicVarieties/ModularityArgument/ThresholdArguments/DenseThresholdCases/CuspUnique.lean}{prime cyclic step} and \href{https://github.com/Deicyde/riemann-hypothesis/blob/40e56bab026552e089d7afcfeff08d2defb9c005/src/Development/AutomorphicArithmetic/AutomorphicVarieties/ModularityArgument/BaseChangeUnits/SolvableEquivalenceOfBaseChange.lean}{solvable degree induction} give the arithmetic formulation in terms of the specified $r$ and its restriction.
\item J.~Arthur and L.~Clozel, \href{https://www.math.toronto.edu/arthur/pdf/30.pdf}{\emph{Simple Algebras, Base Change, and the Advanced Theory of the Trace Formula}}, Annals of Mathematics Studies 120 (1989), Chapter 3, for prime-degree cyclic base change. L.~Clozel and C.~S.~Rajan, \href{https://doi.org/10.1515/crelle-2020-0023}{\emph{Solvable base change}}, Theorem 3, describes the image and fibers for general solvable base change. The precise prime-degree Satake and infinity-type statements used here are in the release's \href{https://github.com/Deicyde/riemann-hypothesis/blob/40e56bab026552e089d7afcfeff08d2defb9c005/src/Development/AutomorphicArithmetic/AutomorphicVarieties/ModularityArgument/BaseChangeUnits/CycCOfCyclic.lean}{cyclic base-change hypotheses}.
\item The release's \href{https://github.com/Deicyde/riemann-hypothesis/blob/40e56bab026552e089d7afcfeff08d2defb9c005/src/Development/AutomorphicArithmetic/AutomorphicVarieties/ModularityArgument/ThresholdArguments/DenseThresholdCases/SVRSol.lean}{polarized attachment at the cyclic descent step} and \href{https://github.com/Deicyde/riemann-hypothesis/blob/40e56bab026552e089d7afcfeff08d2defb9c005/src/Development/AutomorphicArithmetic/AutomorphicVarieties/ModularityArgument/ThresholdArguments/DenseThresholdCases/VWCsdOf.lean}{passage from a polarized Galois representation to a conjugate self-dual cusp form} supply the two different attachment assertions.
\item The \href{https://github.com/Deicyde/riemann-hypothesis/blob/40e56bab026552e089d7afcfeff08d2defb9c005/src/Development/AutomorphicArithmetic/AutomorphicVarieties/ModularityArgument/AutomorphyLifting/WeakOfAtt.lean}{definitions of the flag, power condition and automorphic window}, the \href{https://github.com/Deicyde/riemann-hypothesis/blob/40e56bab026552e089d7afcfeff08d2defb9c005/src/Development/AutomorphicArithmetic/AutomorphicVarieties/ModularityArgument/ThresholdArguments/DenseThresholdCases/TAUpDownOf.lean}{descent of the local window}, the \href{https://github.com/Deicyde/riemann-hypothesis/blob/40e56bab026552e089d7afcfeff08d2defb9c005/src/Development/AutomorphicArithmetic/AutomorphicVarieties/ModularityArgument/LeviSmoothLocalization/WinUpDown.lean}{strict uniformizer bounds}, and the \href{https://github.com/Deicyde/riemann-hypothesis/blob/40e56bab026552e089d7afcfeff08d2defb9c005/src/Development/AutomorphicArithmetic/AutomorphicVarieties/ModularityArgument/ThresholdArguments/DenseThresholdCases/DnDown.lean}{final descent of both witnesses} specify the local part of the conclusion.
\end{enumerate}
